\documentclass[11pt]{article}

% ---------------------------------------------------------------
%  Optimal Approximation for Fully Dynamic Matching and
%  Ordered Ruzsa--Szemeredi Graphs:
%  A Feasibility-Checked Research Execution Plan
%
%  Prepared for: Ruifeng Cao
%  Date: June 10, 2026
%
%  Citation policy used in this document:
%    [V]  Verified against the published source / arXiv during the
%         preparation of this plan (June 2026 web verification pass).
%    [M]  Cited from standard knowledge; canonical and very likely
%         correct, but page numbers / exact venues should be
%         re-checked before any submission (Phase 0, Task V.1).
%  Every nontrivial claim in the document carries a status tag:
%    (V) verified-from-source, (P) proven here, (S) proof sketch,
%    (C) conjecture, (X) speculation.
% ---------------------------------------------------------------

\usepackage[margin=1.05in]{geometry}
\usepackage[T1]{fontenc}
\usepackage{amsmath,amssymb,amsthm,mathtools}
\usepackage{booktabs}
\usepackage{longtable}
\usepackage{enumitem}
\usepackage{xcolor}
\usepackage[expansion=false]{microtype}
\usepackage[colorlinks=true,linkcolor=blue!60!black,citecolor=blue!60!black,urlcolor=blue!50!black]{hyperref}

\newtheorem{theorem}{Theorem}[section]
\newtheorem{lemma}[theorem]{Lemma}
\newtheorem{corollary}[theorem]{Corollary}
\newtheorem{proposition}[theorem]{Proposition}
\newtheorem{conjecture}[theorem]{Conjecture}
\newtheorem{hypothesis}[theorem]{Hypothesis}
\theoremstyle{definition}
\newtheorem{definition}[theorem]{Definition}
\newtheorem{problem}[theorem]{Problem}
\newtheorem{task}{Task}
\theoremstyle{remark}
\newtheorem{remark}[theorem]{Remark}
\newtheorem{observation}[theorem]{Observation}

% claim-status tags
\newcommand{\statusV}{\textcolor{green!45!black}{\textsf{[V]}}}
\newcommand{\statusM}{\textcolor{orange!80!black}{\textsf{[M]}}}
\newcommand{\statusP}{\textcolor{blue!60!black}{\textsf{[P]}}}
\newcommand{\statusS}{\textcolor{violet}{\textsf{[S]}}}
\newcommand{\statusC}{\textcolor{red!70!black}{\textsf{[C]}}}
\newcommand{\statusX}{\textcolor{red!70!black}{\textsf{[X]}}}

\newcommand{\ORS}{\mathrm{ORS}}
\newcommand{\RS}{\mathrm{RS}}
\newcommand{\poly}{\mathrm{poly}}
\newcommand{\polylog}{\mathrm{polylog}}
\newcommand{\OMv}{\mathsf{OMv}}
\newcommand{\eps}{\varepsilon}
\newcommand{\Otil}{\widetilde{O}}
\newcommand{\mm}{\mu}

\title{\bf Optimal Approximation for Fully Dynamic Matching and\\
Ordered Ruzsa--Szemer\'edi Graphs:\\[2pt]
\large A Feasibility-Checked Execution Plan for the Deterministic and\\
Lower-Bound Frontiers}
\author{Research plan prepared for Ruifeng Cao}
\date{June 11, 2026 (v1.1; v1.0: June 10)}

\begin{document}
\maketitle

\begin{abstract}
\noindent
Behnezhad and Ghafari (FOCS 2024) and Assadi, Khanna, and Kiss (SODA 2025)
reduced fully dynamic $(1-\eps)$-approximate maximum matching to a purely
combinatorial quantity, the density $\ORS(n)$ of \emph{ordered
Ruzsa--Szemer\'edi graphs}: the current state of the art maintains the
\emph{edges} of a $(1-\eps)$-approximate maximum matching in amortized
$n^{o(1)}\cdot \ORS(n)$ update time, with high probability, against an
adaptive adversary. Two questions remain open: \textbf{(OP-D)} can a
\emph{deterministic} algorithm achieve the same guarantees, and
\textbf{(OP-L)} can current lower-bound techniques (hidden permutations,
RS-graph reductions, $\OMv$-based hardness) be extended to prove polynomial
lower bounds against \emph{deterministic} dynamic algorithms?

This document is an execution plan for attacking both questions. It
contains: (i) a verified map of the 2024--2026 state of the art, including
Pratt's SOSA 2026 proof that ordered and classical RS densities are
polynomially equivalent, which collapses part of the design space;
(ii) precise formalizations of OP-D and OP-L, including a correction of a
common misreading of the SODA 2025 result (it improved the \emph{update
time}, not the approximation ratio); (iii) several original structural
observations, among them a fully proved $\Theta(n^2)$-vs-$\Otil(n)$
deterministic/randomized separation for matching-size estimation in the
adjacency-probe model, and a derived \emph{architectural barrier} showing
that the entire randomized polylog-update pipeline
(BKSW\,/\,Behnezhad\,/\,ABR\,/\,BG\,/\,AKK) has no deterministic drop-in
replacement at its sublinear-oracle layer; (iv) an honest taxonomy of what
kinds of deterministic lower bounds are achievable today versus blocked by
known cell-probe barriers; and (v) four concrete workstreams with tasks,
decision gates, fallback publishable units, a 9--12 month timeline, a risk
register, and a claim-status appendix. All citations were re-verified
against arXiv/DOI metadata in June 2026 except where explicitly tagged.
\end{abstract}

\tableofcontents

% ===============================================================
\section{Scope, reading of the open problem, and one correction}
\label{sec:scope}
% ===============================================================

\subsection{The open problem as posed}

The prompt poses, in the context of ``optimal approximation for fully
dynamic matching and ordered Ruzsa--Szemer\'edi (ORS) graphs'':

\begin{enumerate}[label=(\alph*),itemsep=2pt]
\item \emph{Can deterministic algorithms achieve the same approximation
guarantees} as the randomized ORS-based algorithms of Behnezhad--Ghafari
\cite{BG24} and Assadi--Khanna--Kiss \cite{AKK25}?
\item \emph{Can current lower-bound techniques (hidden permutations, RS
graphs) be extended to prove polynomial communication lower bounds for
deterministic algorithms in the dynamic setting?}
\end{enumerate}

We formalize these as Problems \ref{prob:OPD} (OP-D) and \ref{prob:OPL}
(OP-L) in Section~\ref{sec:problems}, after fixing the model zoo in
Section~\ref{sec:sota}; the distinctions between maintaining
\emph{edges}, maintaining a \emph{size estimate}, and maintaining a
\emph{witness} turn out to be essential to both questions, and several
natural readings of (a) and (b) are either already resolved or provably
hopeless with current technology, so the formalization step is not
bureaucratic.

\subsection{A correction to the premise \statusV}
\label{sec:correction}

The prompt states that Assadi--Khanna--Kiss (SODA 2025) ``used ordered RS
graphs to significantly improve the \emph{approximation ratio} of fully
dynamic matching.'' The verified record is slightly different, and the
difference matters for what OP-D should ask:

\begin{itemize}[itemsep=2pt]
\item Behnezhad--Ghafari \cite{BG24} (FOCS 2024, pp.~314--327;
arXiv:2404.06069) already achieve approximation ratio $1-\eps$ for every
fixed $\eps>0$. They introduce $\ORS$ graphs and maintain the
\emph{edges} of a $(1-\eps)$-approximate maximum matching in amortized
update time
$O\!\big(\sqrt{\,n^{1+O(\eps)}\cdot \ORS(n)\,}\big)$,
with high probability, against an adaptive adversary.
\item Assadi--Khanna--Kiss \cite{AKK25} (SODA 2025, pp.~2971--2990;
arXiv:2406.13573) keep the same approximation ratio $1-\eps$ and the same
adaptive-adversary guarantee but improve the \emph{update time} to
amortized $n^{o(1)}\cdot \ORS(n)$, thereby reducing the dynamic problem
\emph{entirely} to bounding the combinatorial quantity $\ORS(n)$.
\end{itemize}

So the ``optimal approximation'' in the title of this research direction
is $1-\eps$, which both papers already attain (randomized); the open
algorithmic frontier is (i) the update time as a function of $\ORS(n)$,
(ii) the truth about $\ORS(n)$ itself, and (iii) determinism. OP-D below
is phrased accordingly.

\subsection{What changed in 2025--2026 and why it reshapes the plan \statusV}
\label{sec:newdev}

Two verified developments postdate the original formulation of this open
problem and materially change the attack surface:

\begin{enumerate}[itemsep=3pt]
\item \textbf{Pratt (SOSA 2026): ordered $\approx$ unordered.}
Pratt \cite{Pra26} (arXiv:2502.02455, DOI 10.1137/1.9781611978964.27)
proves that if $\ORS(n)\ge \Omega(n^{c})$ for a fixed $c>0$, then for
every fixed $\delta>0$ the \emph{classical} RS density at slightly smaller
(still linear) matching size satisfies
$\RS\big(n,\Theta_{\delta}(n)\big)\ge \Omega(n^{c(1-\delta)})$.
Consequently the ordering relaxation does \emph{not} buy polynomial
density, the question ``is $\ORS(n)=n^{o(1)}$ or $n^{1-o(1)}$?'' from
\cite{BG24} reduces (in the polynomial regime) to the decades-old
Ruzsa--Szemer\'edi question, and the AKK update time is, up to
$n^{o(1)}$, governed by \emph{classical} RS density. For this plan, the
consequence is that Workstream B (Section~\ref{sec:wsB}) should target RS
and ORS jointly, and that ``find a denser \emph{ordered} construction''
is no longer a promising shortcut to faster algorithms. (Exact parameter
translation: Phase-0 Task~\ref{task:verify}.)
\item \textbf{Bernstein--Bhattacharya--Kiss--Saranurak (STOC 2025):
deterministic maximal matching in sublinear time.} \cite{BBKS25}
(arXiv:2504.20780, DOI 10.1145/3717823.3718153) give a deterministic
fully dynamic \emph{maximal} matching in amortized $\Otil(n^{8/9})$
update time, breaking the long-standing $O(n)$ deterministic barrier on
dense graphs, plus a randomized adaptive-adversary algorithm in
$\Otil(n^{3/4})$. This confirms that (i) the deterministic frontier is
active and the community considers it major, and (ii) even
\emph{maximal} matching --- a $1/2$-approximation primitive ---
currently costs deterministic $n^{8/9}$, which calibrates how far OP-D
is from resolution and supplies the natural first stepping stone
(target D1 below).
\end{enumerate}

\subsection{Deliverables of this plan}

The plan is organized as: verified state of the art
(Section~\ref{sec:sota}); formal problems, original observations, and a
proved probe-model separation with its barrier corollary
(Section~\ref{sec:problems}); four workstreams WS-A--WS-D with concrete
tasks (Sections~\ref{sec:wsA}--\ref{sec:wsD}); timeline, decision gates,
and fallback publishable units (Section~\ref{sec:timeline}); risk
register (Section~\ref{sec:risks}); claim-status and notation appendices
(Appendices~\ref{app:claims}--\ref{app:notation}).

% ===============================================================
\section{Verified state of the art (June 2026)}
\label{sec:sota}
% ===============================================================

\subsection{Model zoo: what exactly is being maintained}
\label{sec:models}

Let $G=(V,E)$, $|V|=n$, $|E|=m$, undergo edge insertions and deletions;
$\mm(G)$ denotes the maximum matching size. The literature mixes four
distinct objectives, and conflating them produces false ``contradictions''
between upper and lower bounds. We fix names:

\begin{description}[itemsep=2pt]
\item[\textsc{Edges}$(\eps)$:] maintain an explicit edge set $M\subseteq E$
that is a matching with $|M|\ge (1-\eps)\mm(G)$ after every update.
\item[\textsc{Size}$(\alpha)$ / \textsc{Size}$(\alpha,\beta n)$:] maintain
a number $\hat{\mm}$ with $\alpha\,\mm(G)-\beta n \le \hat{\mm}\le \mm(G)$
(multiplicative $\alpha$, optional additive slack $\beta n$).
\item[\textsc{Frac}$(\alpha)$:] maintain a fractional matching of value at
least $\alpha\,\mm(G)$ (or within $\alpha$ of the fractional optimum).
\item[\textsc{Witness-VC}$(\alpha)$:] maintain an explicit vertex set that
is an $\alpha$-approximate vertex cover.
\end{description}

Orthogonal axes: \emph{amortized vs.\ worst-case} update time;
\emph{oblivious vs.\ adaptive} adversary; \emph{randomized (whp) vs.\
deterministic}. One important calibration~\statusV: the ORS-based
algorithms \cite{BG24,AKK25} are already robust to \emph{adaptive}
adversaries, so OP-D is a pure determinism question --- the usual
intermediate target ``upgrade oblivious to adaptive'' is unavailable as a
stepping stone; the remaining gap between \cite{AKK25} and OP-D is
exactly ``whp $\to$ always-correct, no coins.''

\subsection{Ordered Ruzsa--Szemer\'edi graphs \statusV}
\label{sec:ors}

\begin{definition}[RS and ORS graphs; \cite{RS78,BG24}]
\label{def:ors}
A graph $G$ on $n$ vertices is an \emph{$\RS$ graph with parameters
$(r,t)$} if its edge set can be partitioned into $t$ edge-disjoint
\emph{induced} matchings $M_1,\dots,M_t$, each of size $r$. It is an
\emph{ordered RS ($\ORS$) graph with parameters $(r,t)$} if there is an
ordering $M_1,\dots,M_t$ of $t$ edge-disjoint matchings of size $r$ such
that every $M_i$ is an induced matching in the subgraph
$M_i\cup M_{i+1}\cup\dots\cup M_t$. Write $\RS(n,r)$, $\ORS(n,r)$ for the
maximum such $t$, and $\RS(n)$, $\ORS(n)$ for the regime $r=\Theta(n)$
(constant locality hidden in the $\Theta$, as in \cite{BG24}).
\end{definition}

Every RS graph is an ORS graph under any ordering, so
$\ORS(n,r)\ge\RS(n,r)$. Verified facts:

\begin{itemize}[itemsep=2pt]
\item \emph{Lower bound (constructions).} $\RS(\,\cdot\,,\Theta(n))\ge
n^{1/\Theta(\log\log n)}$, via Fischer, Lehman, Newman, Raskhodnikova,
Rubinfeld, and Samorodnitsky \cite{FLNRRS02} (attribution per \cite{BG24});
this is super-polylogarithmic but $n^{o(1)}$. Behrend-set-based
constructions \cite{Beh46,RS78} give $n^{1-o(1)}$ induced matchings of
size $n^{1-o(1)}$ (slightly sublinear $r$). \statusV\ (regime details:
verify exact constants in Task~\ref{task:verify}).
\item \emph{Upper bounds.} For $r=\Theta(n)$ no upper bound of the form
$n^{1-\Omega(1)}$ is known: whether $\ORS(n)=n^{o(1)}$ or
$\ORS(n)=n^{1-o(1)}$ is wide open \cite{BG24}. For $r$ very close to
$n/2$ (e.g.\ $r\ge(1/4+\delta)n$ and stronger regimes) classical
triangle-removal/regularity arguments force $t$ to be small;
the precise $(r,t)$ frontier is scattered across the literature and is
collected as Task~\ref{task:rs-table}. \statusM\ for the near-$n/2$
regime.
\item \emph{Equivalence.} Pratt \cite{Pra26}: polynomial $\ORS$ density
implies polynomial $\RS$ density at linear matching size
(Section~\ref{sec:newdev}); so in the polynomial regime, ordered and
unordered densities coincide up to $n^{\delta}$ losses. \statusV
\end{itemize}

\subsection{Upper bounds for fully dynamic matching: verified table}
\label{sec:ub-table}

Throughout, ``det'' = deterministic, ``rand'' = randomized (whp),
``adp'' = correct against adaptive adversaries, ``obl'' = oblivious
adversary assumed. All rows verified against the cited source unless
tagged \statusM.

\medskip
\noindent\textbf{(a) \textsc{Edges}$(\eps)$, i.e.\ $(1-\eps)$-approximate
matching maintained explicitly.}
\begin{center}
\small
\begin{tabular}{@{}llll@{}}
\toprule
Update time (amortized unless noted) & Type & Source & Notes\\
\midrule
$O(\sqrt{m}\,\eps^{-2})$ (worst-case) & det & \cite{GP13} & FOCS'13, pp.~548--557\\
$n/(\log^* n)^{\Omega(1)}$ & rand & \cite{ABKL23} & regularity-based \statusM\ (pages)\\
$n/2^{\Omega(\sqrt{\log n})}$ (bipartite) & rand & \cite{Liu24}+\cite{LW17} & via online matrix-vector\\
$O(\sqrt{n^{1+O(\eps)}\ORS(n)})$ & rand, adp & \cite{BG24} & FOCS'24, pp.~314--327\\
$n^{o(1)}\cdot\ORS(n)$ & rand, adp & \cite{AKK25} & SODA'25, pp.~2971--2990\\
\bottomrule
\end{tabular}
\end{center}

\noindent Calibration \statusV: for approximation ratios below $2/3$
(equivalently $\eps>1/3$), update time $n^{1/2+o(1)}$-type bounds follow
from Bernstein--Stein \cite{BS16} ($3/2+\eps$ in $O(m^{1/4})$); for
$\eps\le 1/3$ the best unconditional bound not parameterized by
$\ORS$ is the $n/2^{\Omega(\sqrt{\log n})}$ of \cite{Liu24} (this
calibration is stated in \cite{Pra26}).

\medskip
\noindent\textbf{(b) Maximal matching (a canonical $1/2$-approximation
primitive).}
\begin{center}
\small
\begin{tabular}{@{}llll@{}}
\toprule
Update time & Type & Source & Notes\\
\midrule
$O(\sqrt{m})$ (worst-case) & det & \cite{NS13} & arXiv:1207.1277\\
$\Otil(n^{8/9})$ & det & \cite{BBKS25} & STOC'25; first $o(n)$ det on dense $G$\\
$O(\log n)$ & rand, obl & \cite{BGS11} & SICOMP'18 corrected version \statusM\\
$O(1)$ & rand, obl & \cite{Sol16} & FOCS'16, pp.~325--334\\
$\Otil(n^{3/4})$ & rand, adp & \cite{BBKS25} & \\
\bottomrule
\end{tabular}
\end{center}

\medskip
\noindent\textbf{(c) \textsc{Size} estimation and better-than-$1/2$
ratios in $\polylog$ time (all randomized).}
\begin{center}
\small
\begin{tabular}{@{}llll@{}}
\toprule
Guarantee & Update time & Source & Notes\\
\midrule
\textsc{Size}$(1/2+\Omega(1))$ & $\polylog$ & \cite{Beh23} & SODA'23, pp.~129--162\\
\textsc{Size}$(1+1/\sqrt2-\eps)^{-1}$ bip., $1.973^{-1}$ gen. & $\polylog$ wc, adp & \cite{BKSW23} & SODA'23 + JACM\\
\textsc{Size}$(2-\sqrt2-\eps)$ general & $\polylog$ & \cite{ABR24} & SODA'24, pp.~3040--3061\\
\textsc{Size}$(1-\eps)$ & $m^{1/2-\Omega_{\eps}(1)}$ & \cite{BKS23} & FOCS'23, arXiv:2302.05030\\
\bottomrule
\end{tabular}
\end{center}
Here $2-\sqrt2\approx 0.5857$; \cite{ABR24} also push to
$(2-\sqrt2-\eps)$ for any fixed $\eps>0$ in $\polylog$ time. \statusV

\medskip
\noindent\textbf{(d) Deterministic landscape beyond maximal.}
\begin{center}
\small
\begin{tabular}{@{}llll@{}}
\toprule
Guarantee & Update time & Source & Notes\\
\midrule
\textsc{Edges}$(1/2-\eps)$ \emph{(i.e.\ $(2{+}\eps)$-apx)} & $\poly(\log n,1/\eps)$ & \cite{BHN16} & STOC'16, arXiv:1604.05765\\
better-than-$2$ apx, bipartite & $O(n^{\gamma})$, any $\gamma>0$ & \cite{BHN16} & second result of the paper\\
$(3/2+\eps)$-apx & $O(m^{1/4}\eps^{-2.5})$ & \cite{BS16} & SODA'16, pp.~692--711\\
$(2+\eps)$-apx \textsc{Frac}+VC & $O(1)$ & \cite{BCH17} & IPCO'17\\
det rounding of dyn.\ fractional & $\polylog$ overhead & \cite{BK21} & ICALP'21 \statusM\\
EDCS maintenance, worst-case & $O(m^{1/4})$-type & \cite{GSSU22} & SOSA'22, pp.~12--23 \statusM\\
\bottomrule
\end{tabular}
\end{center}

\noindent\textbf{Two verified summary points.} (i) No deterministic
algorithm is known with $\polylog$ (or even $n^{o(1)}$) update time for
\emph{any} constant-factor guarantee better than $(2+\eps)$, in any of
the four objectives; the deterministic better-than-2 frontier is
$O(n^{\gamma})$/$O(m^{1/4})$-type \cite{BHN16,BS16}. (ii) No
deterministic algorithm with update time $o(\sqrt m)$ is known for
\textsc{Edges}$(\eps)$ at small $\eps$; even matching Liu's randomized
$n/2^{\Omega(\sqrt{\log n})}$ deterministically is open, because the
underlying $\OMv$ algorithm of Larsen--Williams \cite{LW17} (SODA'17,
pp.~2182--2189) is randomized (determinism status of subpolynomial-saving
$\OMv$: verify in Task~\ref{task:verify}). \statusV/\statusM

\subsection{Randomness audit of the ORS pipeline \statusV}
\label{sec:audit}

Where exactly do coins enter \cite{BG24,AKK25} and the $\polylog$
\textsc{Size} algorithms? This audit drives Workstream A.

\paragraph{Update (v1.1).} Task A0's read-through passes 1--2 are
complete; the machine-readable audit table (v0.3, lemma-level
locations, auto-generated by \texttt{audit\_lab}) is reproduced in
Appendix~\ref{app:audit}, with the full technical discussion in the
companion A0 note. Three sharpenings of the prose audit below
\statusV{}: (i) BG24's coins live \emph{only} in its Lemma~7 sampling
plus the estimator/sparsification interfaces --- the ORS charging
(Lemma~11), certificate bookkeeping, and McGregor boosting
(derandomized by Tirodkar, per AKK25 Prop.~2.2) are deterministic;
(ii) the derandomization problem therefore reduces to exactly three
modules (opportunistic sampler; Beh21 estimator interface; vertex
sparsification); (iii) ABR24's Theorem~2 is a \emph{deterministic}
two-pass streaming $(2-\sqrt2-\eps)$-approximation, so the
combinatorial core is deterministic under pass access and the dynamic
coins are forced by the probe access model alone.

\begin{enumerate}[itemsep=2pt]
\item \emph{Sublinear-time estimators as inner oracles.} The
$\polylog$-update pipelines \cite{BKSW23,Beh23,ABR24} repeatedly invoke
sublinear-time matching/size estimators --- most prominently random
greedy maximal matching (RGMM) estimation in the style of Behnezhad
\cite{Beh21} (FOCS'21, pp.~873--884) --- on implicitly represented
(kernel/contracted/peeled) graphs accessed through adjacency probes.
These estimators are Monte Carlo: per call they probe a sublinear number
of adjacency entries and succeed whp.
\item \emph{Random ranks/permutations.} RGMM correctness and locality
analyses are driven by a uniformly random permutation of edges/vertices;
adaptive-adversary robustness in \cite{BKSW23,AKK25} is obtained by
re-randomizing and by outputting only low-sensitivity quantities (sizes),
plus standard amplification.
\item \emph{Sampling for sparsification/peeling.} \cite{BG24,AKK25} use
random vertex/edge sampling inside the ORS-driven phase structure to
bound the number of ``epochs'' by $\ORS$ density; failure probabilities
are union-bounded over epochs.
\end{enumerate}

Consequence (formalized as Corollary~\ref{cor:barrier}): the pipeline's
randomness is concentrated in a clean interface --- ``estimate
$\mm$ of an implicitly represented graph from few probes'' --- and that
interface is \emph{provably} impossible to implement deterministically
in sublinear probes. Derandomizing the pipeline therefore cannot be a
local substitution; it must change the architecture.

\subsection{Lower-bound landscape (verified) and where determinism
currently plays no role}
\label{sec:lb-landscape}

\begin{itemize}[itemsep=3pt]
\item \emph{$\OMv$-conditional.} Henzinger--Krinninger--Nanongkai--%
Saranurak \cite{HKNS15} (STOC'15, pp.~21--30): under the $\OMv$
conjecture, many exact dynamic problems need $n^{1-o(1)}$ update time.
Liu \cite{Liu24} (FOCS'24, pp.~228--243) proves \emph{equivalences}:
maintaining a $(1+\eps)$-approximate \textsc{Witness-VC} in amortized
$n^{1-c}\eps^{-C}$ time is possible iff $\OMv$ is false; maintaining
$(1-\eps)$-approximate bipartite matching in $n^{1-c}$ iff a nontrivial
``dynamic approximate $\OMv$'' algorithm exists. By K\"onig duality the
\emph{value} of the minimum vertex cover equals $\mm$ in bipartite
graphs, and value maintenance \emph{is} sublinearly possible
\cite{BKS23}; the hardness is for \emph{witness} maintenance --- one
more reason the model zoo of Section~\ref{sec:models} matters. \statusV\
(exact quantifiers: Task~\ref{task:verify}.)
\item \emph{Static sublinear (query) lower bounds, randomized.}
Behnezhad--Roghani--Rubinstein \cite{BRR24} (STOC'24, arXiv:2406.08595):
for every $\delta>0$ there is $\eps>0$ such that estimating $\mm$ within
additive $\eps n$ requires $\Omega(n^{2-\delta})$ adjacency-list
queries; near-matching upper bounds in $n^{2-\Omega_{\eps}(1)}$
\cite{BKS23,BRRS23}. These already hold for randomized algorithms, so
they cannot, by themselves, separate det from rand. \statusV
\item \emph{Hidden permutations (the technique named in the prompt).}
Assadi--Sundaresan \cite{AS23} (FOCS'23, pp.~909--932; arXiv:2310.05728)
build \emph{permutation-hiding} graphs directly on RS graphs and prove
that $(1-\eps)$-approximating bipartite matching in semi-streaming needs
$\Omega\!\big(\log(1/\eps)/\log(1/\beta)\big)$ passes, where
$n^{\beta}$ is the RS density available to the construction. This is a
\emph{distributional} (hence randomized-robust) communication argument
in a \emph{bounded-space} model. \statusV
\item \emph{Dynamic streams.} Assadi--Khanna--Li--Yaroslavtsev
\cite{AKLY16} show $n^{2-3\eps}$-space lower bounds for $n^{\eps}$%
-approximation in dynamic streams --- again randomized-robust, and so
strong already for randomized algorithms that there is little
deterministic/randomized daylight left in that model. \statusM
\item \emph{Cell-probe ceilings.} Unconditional cell-probe lower bounds
for dynamic data structures top out around
$\Omega((\log n/\log\log n)^2)$ (Larsen \cite{Lar12}); polynomial
unconditional bounds are a notorious open frontier (P\u{a}tra\c{s}cu's
multiphase program \cite{Pat10}, nondeterministic barriers
\cite{CGL15}). Nothing here exploits determinism, and no known route
gives polynomial \emph{deterministic-only} cell-probe bounds either.
\statusM
\end{itemize}

\noindent\textbf{Reading for OP-L.} Every polynomial lower bound in this
landscape is either conditional ($\OMv$-style), or lives in a restricted
information model (streams, communication, query/probe), and \emph{all}
of them are randomized-robust. There is, verifiably, no published
polynomial lower bound anywhere in the dynamic-matching literature that
applies to deterministic algorithms but not to randomized ones. That gap
is exactly the opportunity --- and Section~\ref{sec:problems} shows it
cannot be closed by ``extending'' the existing techniques unchanged.

% ===============================================================
\section{Formal problems and original structural observations}
\label{sec:problems}
% ===============================================================

\subsection{The two problems, formalized}

\begin{problem}[OP-D: determinism]\label{prob:OPD}
Design a \emph{deterministic} fully dynamic algorithm that maintains the
edges of a $(1-\eps)$-approximate maximum matching with amortized update
time $f(\eps)\cdot \ORS(n)\cdot n^{o(1)}$, i.e.\ matching \cite{AKK25} up
to $n^{o(1)}$ factors --- or refute the possibility in a meaningful
restricted model. Graded sub-targets, each independently publishable:
\begin{description}[itemsep=2pt]
\item[D1 (maximal/half).] Deterministic fully dynamic \emph{maximal}
matching in $n^{o(1)}$ (ultimately $\polylog$) update time; intermediate:
beat the $\Otil(n^{8/9})$ of \cite{BBKS25}.
\item[D2 (fractional).] Deterministic $(1+\eps)$-approximate
\emph{fractional} matching in $\polylog$ update time (current det:\
$(2+\eps)$ in $O(1)$ \cite{BCH17}).
\item[D3 (size, better-than-2).] Derandomize the $\polylog$-time
\textsc{Size}$(1/2+\Omega(1))$ algorithms \cite{Beh23,BKSW23,ABR24} ---
any constant ratio $>1/2$ in deterministic $n^{o(1)}$ time.
\item[D4 (ORS-parameterized).] The full OP-D statement above.
\end{description}
\end{problem}

\begin{problem}[OP-L: deterministic lower bounds]\label{prob:OPL}
Prove super-polylogarithmic (ideally polynomial) lower bounds on the
update/query complexity of \emph{deterministic} fully dynamic matching
algorithms, in a model strong enough to be meaningful, by extending or
replacing the hidden-permutation/RS-reduction/$\OMv$ toolkits. Taxonomy
of admissible answer types (used throughout WS-C):
\begin{description}[itemsep=2pt]
\item[L1.] Unconditional deterministic lower bounds in oracle/probe
models (achievable now; see Lemma~\ref{lem:probe}).
\item[L2.] Architectural barriers: theorems that the \emph{existing
algorithmic frameworks} cannot be derandomized (achievable now; see
Corollary~\ref{cor:barrier}).
\item[L3.] Conditional bounds under a new deterministic hardness
hypothesis (``Det-$\OMv$''), with deterministic reductions.
\item[L4.] $\ORS/\RS$ \emph{density lower bounds} (constructions): these
lower-bound the runtime of the leading algorithm family
\cite{BG24,AKK25}, not of all algorithms --- a legitimate but weaker
currency.
\item[L5.] Extensions of permutation hiding \cite{AS23} to dynamic
settings with \emph{bounded space or bounded recourse}; we prove below
why some such restriction is forced.
\end{description}
\end{problem}

\subsection{Why the existing techniques cannot be ``extended unchanged'':
a no-go observation \statusP}

\begin{observation}[Randomized-robust techniques cannot give
deterministic-only polynomial bounds at ratios below $2-\sqrt2$]
\label{obs:nogo}
Call a lower-bound technique \emph{randomized-robust} if any bound it
proves for deterministic algorithms it also proves (up to constants) for
randomized ones --- as is the case for distributional/communication
arguments via Yao's principle (hidden permutations \cite{AS23},
RS-reduction streaming bounds, \cite{BRR24}) and for $\OMv$-style
reductions, whose source hypothesis is itself about randomized
algorithms \cite{HKNS15}. Then no randomized-robust technique can prove
a polynomial update-time lower bound for deterministic
\textsc{Size}$(2-\sqrt2-\delta)$ maintenance for any $\delta>0$:
such a bound would equally apply to randomized algorithms and contradict
the verified $\polylog$-time randomized algorithm of \cite{ABR24}.
The same argument with \cite{BKSW23,Beh23} blocks all ratios in
$(1/2,\,2-\sqrt2\,]$ and with \cite{BBKS25} blocks
$\omega(n^{3/4+o(1)})$ bounds for adaptive randomized maximal matching.
\end{observation}

\noindent\emph{Consequences.} (i) A deterministic polynomial lower bound
at ratio better than $1/2+\delta$ must use determinism \emph{essentially}
--- e.g.\ adversary arguments that adapt to the algorithm's fixed
behaviour, fooling-set/rank-style deterministic communication bounds, or
a deterministic hardness hypothesis (L3). (ii) For ratios in
$(2-\sqrt2,\,1)$, polynomial lower bounds are open \emph{even for
randomized} algorithms (Liu \cite{Liu24} covers the $\eps\to0$ end only
conditionally/by equivalence), so OP-L at those ratios subsumes a major
open randomized question; the honest deterministic-specific opportunity
is at ratios where randomized algorithms are already fast, i.e.\
$[1/2,\,2-\sqrt2]$, and in the parameterized $\ORS$ regime.

\subsection{A composition observation linking ORS to OMv \statusS}

\begin{observation}[If ORS is sparse, dynamic approximate $\OMv$ is easy]
\label{obs:compose}
Suppose $\ORS(n)=n^{o(1)}$. Then \cite{AKK25} maintains \textsc{Edges}%
$(\eps)$ in amortized $n^{o(1)}$ time; feeding this into Liu's
equivalence \cite{Liu24} (algorithm direction) yields a nontrivial
``dynamic approximate $\OMv$'' algorithm in the sense of \cite{Liu24}.
Contrapositive: hardness of dynamic approximate $\OMv$ would imply
$\ORS(n)=n^{\Omega(1)}$, i.e.\ a \emph{conditional combinatorial lower
bound} for RS-type densities --- an unusual direction (complexity
$\Rightarrow$ extremal combinatorics) worth a short note if the
parameters compose. Status: sketch; the quantifier matching
($\eps$ dependence, value-vs-edges, bipartiteness) is exactly Phase-0
Task~\ref{task:verify}(c).
\end{observation}

\subsection{A proved deterministic/randomized separation in the probe
model, and the architectural barrier \statusP}
\label{sec:probelemma}

The following is elementary but, to our knowledge after the June 2026
verification pass, not stated in this form in the dynamic-matching
literature; its role here is to convert the randomness audit of
Section~\ref{sec:audit} into a theorem-grade barrier and to seed
Workstream C. (Closest precedent: deterministic-vs-randomized gaps for
\emph{learning} a hidden matching \cite{ABKRS04}, a harder
reconstruction problem.)

\begin{definition}[Adjacency-probe model]
An algorithm accesses an unknown graph $G$ on vertex set $[n]$ only via
probes $\textsc{Adj}(u,v)\in\{0,1\}$, chosen adaptively. Cost $=$ number
of probes.
\end{definition}

\begin{lemma}[Quadratic deterministic probe bound for gap matching]
\label{lem:probe}
Any deterministic adaptive algorithm that distinguishes
$\mm(G)=0$ from $\mm(G)\ge n/4$ must make more than
$\binom{\lceil n/2\rceil}{2}-1 \;=\;\Omega(n^2)$ probes in the worst
case. Consequently no deterministic algorithm with $o(n^2)$ probes can
estimate $\mm(G)$ within additive $n/8$, nor within any multiplicative
factor with additive slack $<n/8$. In contrast, the same gap problem is
solvable by a randomized algorithm with $O(n\log n)$ probes (error
probability $1/\poly(n)$), and $\Omega(n)$ probes are necessary for
randomized algorithms; so the deterministic and randomized probe
complexities of the gap problem are $\Theta(n^2)$ and
$\widetilde{\Theta}(n)$ respectively --- a quadratic separation.
\end{lemma}

\begin{proof}
\emph{Deterministic lower bound.} Run the deterministic algorithm
against the adversary that answers $0$ to every probe. Since the
algorithm is deterministic, the probe sequence and final output are
determined; let $P$ be the set of probed pairs, and suppose
$|P|\le\binom{\lceil n/2\rceil}{2}-1$. The empty graph $G_0$ is
consistent with all answers and has $\mm(G_0)=0$. We claim the graph
$H=K_n\setminus P$ (all unprobed pairs present) is also consistent ---
trivially, every probed pair is a non-edge of $H$ --- and satisfies
$\mm(H)\ge n/4$. Suppose not: $\mm(H)<n/4$. Taking both endpoints of a
maximum matching of $H$ yields a vertex cover $C$ of $H$ with
$|C|\le 2\mm(H)<n/2$. Then $U=V\setminus C$ has $|U|>n/2$, hence
$|U|\ge\lceil n/2\rceil$ \emph{(integrality)}. No edge of $H$ lies
inside $U$ (it would be uncovered), so \emph{every} pair inside $U$
belongs to $P$, giving
$|P|\ge\binom{|U|}{2}\ge\binom{\lceil n/2\rceil}{2}$, a contradiction.
So both $G_0$ ($\mm=0$) and $H$ ($\mm\ge n/4$) are consistent with the
transcript, and whatever the algorithm outputs is wrong for one of them.

\emph{Randomized upper bound.} Under the promise, if $\mm(G)\ge n/4$
then $G$ contains $\ge n/4$ pairwise-disjoint edges, so a uniformly
random pair is an edge with probability at least
$(n/4)/\binom{n}{2}\ge 1/(2n)$. Probing $cn\log n$ independent uniform
pairs finds an edge with probability $1-n^{-\Omega(c)}$ if
$\mm(G)\ge n/4$, and finds none if $\mm(G)=0$; output accordingly.
(Distinguishing only requires certifying $\mm\ge1$ under the promise.)

\emph{Randomized lower bound.} By Yao's principle take the distribution:
with probability $1/2$ the empty graph; with probability $1/2$ a
uniformly random matching of size $n/4$ (on a uniformly random
$n/2$-subset). Any fixed probe is an edge with probability at most
$(n/4)\cdot\frac{1}{\binom{n}{2}}\cdot\Theta(1)=O(1/n)$ under the second
component; with $o(n)$ probes the transcript is all-zeros with
probability $1-o(1)$ under \emph{both} components, forcing error
$\ge 1/2-o(1)$.
\end{proof}

\begin{remark}
The bound $\binom{\lceil n/2\rceil}{2}\approx n^2/8$ is tight up to the
constant: probing all pairs costs $\binom n2$ and a slightly cleverer
deterministic strategy can certify the gap after
$\binom n2-\binom{\lceil n/2\rceil}{2}+O(n)$ probes
\emph{(verify exact constant; minor)}. The lemma is robust to making
the promise bipartite, to replacing $n/4$ by $\delta n$, and to allowing
the algorithm to output ``a matching'' instead of a bit. \statusS\ for
these routine variants.
\end{remark}

\begin{corollary}[Architectural barrier for derandomizing the
$\polylog$/ORS pipelines]
\label{cor:barrier}
Fix any dynamic-matching algorithm whose update/query procedure invokes,
as a black box, a matching-size gap estimator on an implicitly
represented $N$-vertex graph accessed via adjacency probes --- as do
\cite{BKSW23,Beh23,ABR24,BG24,AKK25} via RGMM-style sublinear estimation
\cite{Beh21} (Section~\ref{sec:audit}). If the black box is replaced by
any \emph{deterministic} probe algorithm with the same interface and
guarantee, then by Lemma~\ref{lem:probe} some calls cost
$\Omega(N^2)$ probes in the worst case, destroying the sublinear
per-call budget that the frameworks' analyses rely on. Hence there is no
deterministic \emph{drop-in} replacement at the oracle layer:
derandomizing these pipelines requires re-architecting them around
\emph{incrementally maintained deterministic summaries} rather than
fresh per-call sublinear estimation. \statusP\ (as a statement about the
probe interface; it is \emph{not} an unconditional lower bound for OP-D,
and we will say so wherever it is used.)
\end{corollary}

\begin{remark}[Pseudorandomness does not contradict the barrier]
A ``deterministic'' algorithm that fixes a PRG seed is still a
deterministic algorithm, and Lemma~\ref{lem:probe} applies: the
worst-case adversary may depend on the fixed seed. This cleanly explains
why naive ``just plug in Nisan/$k$-wise ranks'' plans for RGMM fail at
the model level, while leaving open the legitimate route of
\emph{maintained} structures (WS-A3). \statusP
\end{remark}

\subsection{Consistency checks that calibrate difficulty \statusP}
\label{sec:consistency}

\begin{observation}[D2 implies beating the randomized frontier]
\label{obs:d2hard}
A deterministic $\polylog$-time $(1+\eps)$-approximate \emph{fractional}
matching (target D2), combined with deterministic rounding
\cite{BK21,ACCSW18}, would yield an integral
$\approx\frac{2}{3}(1-O(\eps))$-approximation in $\polylog$ time
(the $3/2$ integrality gap of the matching LP on general graphs is the
bottleneck), i.e.\ ratio $\approx 0.66 > 2-\sqrt2\approx 0.586$ ---
which would improve the best \emph{randomized} $\polylog$ ratio
\cite{ABR24}. Hence D2 is at least as hard as a major open randomized
question; the plan treats D2 as a stretch target and orders work
accordingly. The same calibration shows $\polylog$-time EDCS
maintenance is open even randomized (it would likewise beat
\cite{ABR24}). \statusP\ (modulo the exact rounding loss; verify in
Task~\ref{task:verify}(d).)
\end{observation}

\begin{observation}[What OP-L can and cannot mean, honestly]
\label{obs:honest}
``Polynomial communication lower bounds for deterministic dynamic
algorithms,'' read literally as unconditional polynomial cell-probe
bounds, is beyond current technology regardless of determinism
(Section~\ref{sec:lb-landscape}); and in bounded-space dynamic streams
the randomized bounds \cite{AKLY16} are already near-maximal, leaving no
det/rand daylight. The defensible interpretations are exactly
L1--L5 of Problem~\ref{prob:OPL}; the plan commits to L1+L2
unconditionally (deliverable in months), pursues L3/L5 as research
bets, and uses L4 as the bridge to Workstream B. \statusP\ (as an
assessment, with the cited evidence.)
\end{observation}

% ===============================================================
\section{Workstream A: derandomization (OP-D)}
\label{sec:wsA}
% ===============================================================

Guiding principle from Corollary~\ref{cor:barrier}: do not derandomize
the oracle calls; \emph{remove} them. Every task below replaces
``fresh sublinear estimation'' by ``maintained deterministic structure''
somewhere in the stack.

\begin{task}[A0: randomness audit, formal version]\label{task:audit}
Produce a table mapping every random choice in
\cite{BKSW23,Beh23,ABR24,BG24,AKK25} to (i) the property it buys
(correctness whp / adaptivity / amortization), (ii) the failure mode if
fixed, (iii) whether limited independence suffices in the analysis.
Output: 4--6 page technical note; feeds A1--A4 and is reusable as a
related-work section. Effort: 2--3 weeks. Risk: low.
\end{task}

\begin{task}[A1: deterministic fractional route, graded]\label{task:frac}
Current det.\ frontier: $(2+\eps)$-\textsc{Frac} in $O(1)$ \cite{BCH17};
det.\ rounding exists \cite{BK21}. Subgoals, in order:
(a) det.\ $(2-\delta_0)$-\textsc{Frac} in $\polylog$ for some explicit
$\delta_0>0$ via primal-dual with deterministic level resets
(\cite{BHI15,BHN16} machinery);
(b) understand precisely why the \cite{BHN16} better-than-2 bipartite
algorithm pays $n^{\gamma}$ --- is the bottleneck an estimation
primitive (then Corollary~\ref{cor:barrier} predicts an obstruction) or
amortized rebuilding (then deterministic multi-scale rebuilding may
help)?
(c) only then attempt D2, knowing Observation~\ref{obs:d2hard} prices
it. Effort: 2--3 months for (a)--(b). Risk: medium.
\end{task}

\begin{task}[A2: EDCS route]\label{task:edcs}
Edge-degree-constrained subgraphs give $(3/2+\eps)$ guarantees and are
maintainable deterministically in $O(m^{1/4})$-type worst-case time
\cite{GSSU22,BS16}. Subgoal: deterministic EDCS maintenance in
$n^{1/2-\delta}$, attacking the rebuild bottleneck with deterministic
expander-decomposition-style partitioning (as in \cite{BBKS25}, whose
$n^{8/9}$ comes from exactly such tools). Calibration: full $\polylog$
EDCS is open even randomized (Observation~\ref{obs:d2hard}), so the
target is the exponent, not $\polylog$. Effort: 3 months. Risk:
medium-high.
\end{task}

\begin{task}[A3: the crux --- rank-robust greedy / D1]\label{task:rank}
Deterministic dynamic \emph{maximal} matching in $n^{o(1)}$ (D1) is the
keystone: it is the minimal primitive whose randomized version
($O(1)$ \cite{Sol16}) powers the kernels of the better-than-$1/2$
pipelines, and \cite{BBKS25} just moved the det.\ bound to
$\Otil(n^{8/9})$. Attack plan:
(a) reprove \cite{BBKS25} from our audit template and isolate where
$8/9$ arises (their tradeoff between deterministic expander
decompositions and rebuild schedules);
(b) formalize the obstruction: for any \emph{fixed} edge ranking
$\pi$, a rank-aware adversary forces greedy-along-$\pi$ to suffer
$\Omega(n)$ recourse per update --- prove this as a standalone lower
bound for the natural ``fixed-rank greedy'' algorithm family
\statusS;
(c) circumvent (b) with multi-scale rank hierarchies (Bentley--Saxe-style
deterministic rebuilding trading recourse for a $(1-1/\polylog)$
approximation slack), targeting det.\ \emph{almost}-maximal
($(1/2-o(1))$-approx) in $n^{3/4}$, then $n^{1/2+o(1)}$.
Effort: the core 4--5 months of WS-A. Risk: high, with publishable
intermediates at every exponent improvement and at the lower bound (b).
\end{task}

\begin{task}[A4: limits of pseudorandomness, short note]\label{task:prg}
Determine the minimal independence/seed length under which the RGMM
locality analyses \cite{Beh21} survive. \emph{Verified obstacle
(v1.1)} \statusV{}: Beh21's Lemma~3.10 charges queries via the map
$\phi(\pi,\vec P)$ that rotates ranks along query paths --- a
bijection of the \emph{full} symmetric group; $k$-wise families are
not closed under these rotations, so the $O(\bar d\log n)$ analysis
does not transfer to limited independence as-is, while experiment E3
of \texttt{audit\_lab} finds even $2$-wise ranks empirically
matching full independence on oblivious instances --- a genuine
analysis/experiment gap that this task should close in either
direction. Combine with the
Remark after Corollary~\ref{cor:barrier} into a note ``what
pseudorandomness can and cannot do for dynamic matching.'' Effort: 3--4
weeks, parallelizable. Risk: low (the note is valuable even if the
answer is ``full independence needed'').
\end{task}

\begin{task}[A5: deterministic ORS-parameterized algorithm (D4)]
\label{task:d4}
Only after A3 yields a det.\ kernel primitive: re-derive the
\cite{AKK25} epoch structure with deterministic peeling, replacing
whp-union-bounds over epochs by potential functions. The honest
expectation \statusC{} is det.\ \textsc{Edges}$(\eps)$ in
$\Otil(\sqrt{n\cdot\ORS(n)})$-type time first (the \cite{BG24} shape),
since \cite{AKK25}'s improvement leans hardest on random sampling.
\emph{Refined target (v1.1)} \statusV{}: by the A0 audit, D4 reduces
to derandomizing three modules; the only one not covered by
Lemma~\ref{lem:probe} is \textsc{Det-Opportunistic-Matching} ---
given adjacency access and the \emph{promise} $\mu(G[U])\ge\delta
n$, deterministically output a $\gamma\delta n$-matching in
$\Otil(m\,n^{O(\gamma)}/\Delta_{\mathrm{IN}}(M))$ time (AKK25
Lemma~3.1 shape; $\Otil(n^{2+O(\gamma)}/d_{\mathrm{avg}})$ for the
BG24 Lemma~7 shape). This is a promise-\emph{search} problem, outside
the decision lower bound of Lemma~\ref{lem:probe}; the worst case
$\Omega(m)$ is necessary even randomized \cite{ACK19}, so only the
$\Delta_{\mathrm{IN}}$-adaptivity needs replication.
\end{task}

% ===============================================================
\section{Workstream B: the combinatorics of \texorpdfstring{$\ORS/\RS$}{ORS/RS} density}
\label{sec:wsB}
% ===============================================================

By \cite{AKK25}+\cite{Pra26}, \emph{any} progress on RS density at linear
matching size transfers to the dynamic-matching frontier in both
directions; this workstream is the highest-variance, highest-reward
component, and also produces the experimental deliverables.

\begin{task}[B0: the $(r,t)$ frontier table]\label{task:rs-table}
Compile, with verbatim verified statements, the best known bounds on
$\RS(n,r)$ and $\ORS(n,r)$ across regimes: $r=n/2-o(n)$ (constant $t$
regimes), $r=\delta n$ (the \cite{FLNRRS02} $n^{1/\Theta(\log\log n)}$
construction; the exact $\delta$-dependence used by \cite{AKK25}),
$r=n^{1-o(1)}$ (Behrend regime, $t=n^{1-o(1)}$
\cite{Beh46,RS78}), including corner-free-set improvements
\cite{CFTZ22}. Identify \emph{the} regime that controls the
\cite{AKK25} runtime (their $\eps$-dependence) --- this is the
single most important verification in the plan. Output: survey
section + machine-readable table. Effort: 3--4 weeks. Risk: low.
\end{task}

\begin{task}[B1: \texttt{ors\_lab} --- exact small-$n$ values]
\label{task:orslab}
Build a verified search package computing exact $\ORS_r(n)$ and
$\RS_r(n)$ for $n\lesssim 40$:
(a) two \emph{independent} encodings (SAT via incremental CNF with
symmetry breaking; CP-SAT/ILP) cross-validated on all instances, per
the standing methodology of this research program;
(b) certificates: witness colourings checked by a separate
$O(t\cdot n^2)$ verifier; UNSAT instances logged with solver proofs
where available;
(c) outputs: exact tables, extremal examples, and a targeted search for
small-$n$ gaps $\ORS_r(n)>\RS_r(n)$ --- Pratt's equivalence is
asymptotic, so small-$n$ separations are possible and would inform
constructions (and make a nice figure regardless of sign). \statusC{}
on finding separations. Effort: 4--6 weeks including writeup. Risk: low
(the tables are a contribution per se; nearest analogue: exact
small-$n$ extremal tables in the induced-matchings literature).
\end{task}

\begin{task}[B2: constructions]\label{task:construct}
Lift recent corner-free/cap-set technology \cite{CFTZ22,BCCGNSU17} into
the $r=\delta n$ regime: the concrete question is whether combinatorial
degenerations improve the $n^{1/\Theta(\log\log n)}$ of \cite{FLNRRS02}
for \emph{any} fixed $\delta>0$. Even a $\log$-factor-in-the-exponent
improvement directly worsens the conditional optimality picture for
\cite{AKK25} and is a standalone combinatorics paper. \statusC.
Effort: open-ended; timebox 2 months with B1's extremal examples as
seeds. Risk: high.
\end{task}

\begin{task}[B3: upper bounds where ordering should not help]
\label{task:upper}
The classical upper bounds for $\RS(n,\delta n)$ go through
triangle-removal/regularity; the ordered variant breaks the symmetry
those proofs use, which is exactly why \cite{BG24} had to ask the ORS
question separately --- and \cite{Pra26} now supplies the reduction
ORS$\to$RS, so the leverage point is: prove $\RS(n,\delta n)=n^{o(1)}$
for $\delta$ below the current threshold, or a conditional/regime-%
restricted version, then transfer. Any such bound \emph{upper-bounds the
runtime of \cite{AKK25}} (an L4-type result: it makes the algorithm
unconditionally fast, the optimistic resolution of the program). 
Effort: research bet; timebox 6 weeks initial. Risk: high. \statusC
\end{task}

% ===============================================================
\section{Workstream C: deterministic-specific lower bounds (OP-L)}
\label{sec:wsC}
% ===============================================================

\begin{task}[C1: probe-model separations, paper-grade (L1)]
\label{task:probe}
(a) Write up Lemma~\ref{lem:probe} with the variants (bipartite,
$\delta n$ gaps, output-a-matching), the tightness remark, and the
\cite{ABKRS04} comparison. \emph{Prior-art scoping (v1.1)}
\statusV{}: Assadi--Solomon \cite{AS19} (Thm.~3) already prove a
deterministic $\Omega(n^2)$ vs.\ randomized $\Otil(n\beta)$
separation for \emph{finding} a maximal matching on
adjacency-\emph{array} graphs of neighborhood independence
$\beta=2$, and \cite{ACK19} (Thm.~6) shows finding a maximal
matching needs $\Omega(n^2)$ queries even randomized. The note's
novelty claim must therefore rest on: the adjacency-\emph{matrix}
probe model, the size-\emph{estimation} (decision) task --- where
randomized is $\Otil(n)$ \cite{Beh21}, unlike finding --- the
explicit clique-coverage threshold
$\binom{\lceil n/2\rceil}{2}$, and the pipeline barrier of
Corollary~\ref{cor:barrier}; cite both works prominently.
(b) \emph{List/adjacency-array model:} the matrix-model adversary is too
weak there (degrees reveal emptiness in $O(n)$), so plant the gap on top
of a fixed background graph: target --- any deterministic algorithm
achieving \textsc{Size}$(1/2+\delta)$ needs $n^{1+\Omega(1)}$
(stretch: $\Omega_{\delta}(n^2)$) adjacency-list probes; method ---
derandomized congestion adversaries adapted from the dummy-vertex
machinery of \cite{BRR24}, replacing their distributional core by an
explicit potential the adversary maximizes against the fixed algorithm.
\statusS\
(c) \emph{Exhaustive small-$n$ verification:} implement the adversary
and verify the exact bound $\binom{\lceil n/2\rceil}{2}$ for
$n\le 12$ by exhaustive simulation (ties into WS-D).
Deliverable: a SOSA-style note ``Deterministic probe complexity of
approximate matching size'' = L1+L2 packaged
(Lemma~\ref{lem:probe} + Corollary~\ref{cor:barrier} + (b) if it
lands). Effort: (a),(c) 3--4 weeks; (b) 2 months. Risk: low/medium.
\end{task}

\begin{task}[C2: dynamic permutation hiding under bounded recourse (L5)]
\label{task:hiding}
Observation~\ref{obs:honest} forces a resource restriction. Proposed
model: \emph{bounded-recourse deterministic dynamic algorithms} --- the
algorithm maintains an explicit matching and may change $\le \rho(n)$
edges per update; lower-bound $\rho$ for approximation $1-\eps$ via a
$k$-phase adaptive game in which the adversary inserts the layers of a
permutation-hiding ORS gadget \cite{AS23} in an order chosen
\emph{adaptively against the deterministic transcript} (this is the step
Yao-based proofs cannot do, hence the deterministic-specific hope).
Milestones: (a) formal model + trivial bounds; (b) recourse lower bound
$\rho=\Omega(\text{poly}(\ORS\text{-density}))$ for fixed-rank greedy
(connects to Task~\ref{task:rank}(b)); (c) general deterministic
bounded-recourse bound \statusX. Honest wall, stated upfront: recourse
bounds do not bound update \emph{time} for \textsc{Size} algorithms;
they target \textsc{Edges}. Effort: 3 months. Risk: high.
\end{task}

\begin{task}[C3: a Det-$\OMv$ hypothesis (L3)]\label{task:detomv}
(a) Formulate: ``no deterministic algorithm solves $\OMv$ in
$n^{3-\delta}$ total time (with polynomial preprocessing)'' and the
$\eps$-approximate variants matching \cite{Liu24}; first sanity step is
a focused search for deterministic subcubic(-by-polynomial-factor) OMv
--- note \cite{LW17} is randomized, and the deterministic status of even
$n^3/2^{\Omega(\sqrt{\log n})}$ must be pinned down
(Task~\ref{task:verify}(e)).
(b) Check the \cite{HKNS15,Liu24} reductions are deterministic (they
are reduction-side deterministic on inspection \statusS), so Det-$\OMv$
transfers to deterministic dynamic lower bounds with the same
quantitative shape.
(c) Gadget program: strengthen specific reductions using
\emph{deterministic} communication primitives (Equality/Greater-Than
have $\Theta(n)$ deterministic but $O(\log n)$/$O(1)$ randomized
communication) to get deterministic-only separations in intermediate
communication-game models. \statusC\ Effort: (a),(b) 1 month; (c)
open-ended, timebox 2 months. Risk: medium --- the failure mode
``hypothesis is false'' is itself a (smaller) result.
\end{task}

% ===============================================================
\section{Workstream D: engineering, verification, packaging}
\label{sec:wsD}
% ===============================================================

Per the standing methodology of this research program (every paper ships
runnable, cross-validated code):

\begin{itemize}[itemsep=2pt]
\item \texttt{ors\_lab/}: Task~\ref{task:orslab}; SAT (PySAT/kissat) and
CP-SAT encoders, certificate verifier, table generator, plots.
\item \texttt{probe\_lb/}: Task~\ref{task:probe}(c); adversary
implementation, exhaustive verification harness $n\le 12$, and a fuzz
harness pitting published randomized estimators against the adversary
instances (illustrative figures for the note).
\item \texttt{audit/}: machine-readable randomness-audit table
(Task~\ref{task:audit}) rendered to LaTeX.
\item Reproducibility: pinned environments, seeds, CI on small
instances; independent re-implementation rule for any number that
enters a paper.
\end{itemize}
Effort: woven into B1/C1; ${\sim}3$ weeks dedicated. Risk: low.

% ===============================================================
\section{Phase 0 verification, timeline, gates, fallbacks}
\label{sec:timeline}
% ===============================================================

\begin{task}[V.1 (Phase 0): verification pass --- blocking]
\label{task:verify}
Before any proof is trusted: (a) exact \cite{AKK25} theorem statements
($\eps$-dependence, the precise ORS regime $r(\eps)$, adaptive
guarantee); (b) exact \cite{Pra26} parameter translation; (c) exact
\cite{Liu24} equivalence quantifiers (value vs.\ witness vs.\ edges;
bipartite only?) for Observation~\ref{obs:compose}; (d) exact rounding
losses in \cite{BK21,ACCSW18} for Observation~\ref{obs:d2hard}; (e)
deterministic status of subcubic-savings $\OMv$ (for
Task~\ref{task:detomv}); (f) re-check all \statusM-tagged citations
(page numbers, venues); (g) arXiv sweep for 2025--26 successors of
\cite{BBKS25,ABR24,AKK25}. Output: corrected Section~\ref{sec:sota} +
errata list. Effort: 2--3 weeks.
\end{task}

\subsection{Timeline (9--12 months) and decision gates}

\begin{center}
\small
\begin{tabular}{@{}lll@{}}
\toprule
Phase & Months & Content and exit gate\\
\midrule
M0 & 0--1 & Task~\ref{task:verify}; A0 audit started.
\emph{Gate G0:} premises confirmed; else re-plan.\\
M1 & 1--3 & C1(a,c) note drafted; B1 \texttt{ors\_lab} v1 ($n\le30$);
A0 done. \emph{Gate G1:} if C1(b)\\
 & & resists 4 focused weeks, pivot C1(b)$\to$C2 recourse model.\\
M2 & 3--6 & A3 (a,b): \cite{BBKS25} reproved, fixed-rank recourse LB;
B2 timebox. \emph{Gate G2:}\\
 & & exponent progress on D1? If yes, A3(c) becomes the main paper
track.\\
M3 & 6--9 & A1/A2 in parallel with C3(a,b); B3 timebox.
\emph{Gate G3:} pick main paper.\\
M4 & 9--12 & Integration, writing, code release; submit main +
note(s).\\
\bottomrule
\end{tabular}
\end{center}

\subsection{Fallback publishable units (pre-committed)}
\begin{enumerate}[itemsep=2pt]
\item \emph{SOSA/short:} ``Deterministic probe complexity of approximate
matching size and a barrier for derandomizing dynamic matching''
(Lemma~\ref{lem:probe} + Corollary~\ref{cor:barrier} + variants) ---
deliverable from M1 material alone.
\item \emph{Combinatorics/experimental note:} exact small-$n$
$\ORS/\RS$ tables with certificates and any ordered-vs-unordered
small-$n$ separation (B1).
\item \emph{Position/SoK only if both stall:} ``What would it take to
derandomize dynamic matching?'' --- the audit + no-go observations +
honest taxonomy; lowest value, explicitly last resort.
\end{enumerate}

% ===============================================================
\section{Risk register}
\label{sec:risks}
% ===============================================================

\begin{center}
\small
\begin{tabular}{@{}p{0.30\textwidth}p{0.12\textwidth}p{0.46\textwidth}@{}}
\toprule
Risk & Likelihood / impact & Mitigation\\
\midrule
$\ORS/\RS$ density resolved by others (either direction) & low / very
high & Monthly arXiv monitor (combinatorics + data structures); plan
survives both outcomes: sparse $\Rightarrow$ B3/L4 framing, dense
$\Rightarrow$ constructions paper.\\
C1(b) list-model LB harder than hoped & medium / medium & Pre-planned
pivot at G1 to the recourse model (C2), which reuses the machinery.\\
Det-$\OMv$ false (det.\ subcubic-savings OMv exists) & medium / medium &
Sanity-search first (C3(a)); a refutation is itself a small result and
redirects C3(c).\\
\cite{BBKS25} techniques extended below $n^{8/9}$ during the project &
medium / medium & A3(a) keeps us conversant; exponent race has many
rungs; recourse LB A3(b)/C2 is differentiated.\\
A \statusM{} citation or a premise fails verification & low / high &
Phase 0 is blocking by design; claim-status appendix localizes damage.\\
Scope creep into the randomized-open zone
(ratios $>2-\sqrt2$) & medium / low & Observation~\ref{obs:nogo}
consulted at every gate; such targets require explicit re-scoping.\\
\bottomrule
\end{tabular}
\end{center}

% ===============================================================
\appendix
\section{Randomness audit table (Task A0, v0.3)}
\label{app:audit}
Auto-generated from \texttt{audit\_lab/auditlib/audit\_table.py};
regenerate with \texttt{run\_all.py}. Evidence tags E1--E5 refer to
the experiments of the \texttt{audit\_lab} package; see the A0
technical note for per-row discussion.

\input{audit_table.tex}

\section{Claim-status appendix}
\label{app:claims}
% ===============================================================

Tag semantics: \statusV{} verified against the published source/arXiv in
the June 2026 pass; \statusM{} canonical memory citation, re-check
before submission; \statusP{} proven in this document; \statusS{} proof
sketch, completion is a named task; \statusC{} conjecture/target;
\statusX{} speculation.

\begin{center}
\small
\begin{tabular}{@{}p{0.62\textwidth}l l@{}}
\toprule
Claim & Status & Where\\
\midrule
\cite{BG24}: $(1-\eps)$-\textsc{Edges}, $O(\sqrt{n^{1+O(\eps)}\ORS(n)})$,
whp, adaptive & \statusV & \S\ref{sec:correction}\\
\cite{AKK25}: same guarantee, $n^{o(1)}\ORS(n)$; improvement is the
\emph{update time} & \statusV & \S\ref{sec:correction}\\
\cite{Pra26}: polynomial $\ORS \Rightarrow$ polynomial $\RS$ (linear
$r$) & \statusV & \S\ref{sec:newdev}\\
\cite{BBKS25}: det.\ maximal in $\Otil(n^{8/9})$; rand.\ adaptive
$\Otil(n^{3/4})$ & \statusV & \S\ref{sec:newdev}\\
\cite{ABR24}: \textsc{Size}$(2-\sqrt2-\eps)$ in $\polylog$, general
graphs & \statusV & \S\ref{sec:ub-table}\\
\cite{Liu24}: $\OMv$ equivalences; witness-vs-value distinction &
\statusV & \S\ref{sec:lb-landscape}\\
\cite{AS23}: pass LBs via permutation hiding on RS graphs & \statusV &
\S\ref{sec:lb-landscape}\\
\cite{BRR24}: randomized $\Omega(n^{2-\delta})$ list-queries for
additive $\eps n$ & \statusV & \S\ref{sec:lb-landscape}\\
$(r,t)$ frontier in the near-$n/2$ regime; \cite{ABKL23} pages;
\cite{AKLY16,BGS11,BK21,GSSU22} details; cell-probe refs & \statusM &
Task~\ref{task:verify}(f)\\
Obs.~\ref{obs:nogo} (no-go for randomized-robust techniques) &
\statusP & \S\ref{sec:problems}\\
Lemma~\ref{lem:probe} (quadratic det.\ probe bound; $\Theta(n^2)$ vs
$\widetilde\Theta(n)$) & \statusP & \S\ref{sec:probelemma}\\
Cor.~\ref{cor:barrier} (no det.\ drop-in at the oracle layer;
framework-relative) & \statusP & \S\ref{sec:probelemma}\\
Obs.~\ref{obs:d2hard} (D2 $\Rightarrow$ beating \cite{ABR24}; modulo
exact rounding loss) & \statusP/\statusS & \S\ref{sec:consistency}\\
Obs.~\ref{obs:compose} (ORS sparse $\Rightarrow$ dyn.\ approx.\
$\OMv$) & \statusS & Task~\ref{task:verify}(c)\\
Fixed-rank greedy recourse LB; list-model det.\ LB & \statusS &
Tasks~\ref{task:rank}(b), \ref{task:probe}(b)\\
Small-$n$ $\ORS>\RS$ separations; B2/B3 outcomes; C2(c); C3(c) &
\statusC/\statusX & WS-B, WS-C\\
\bottomrule
\end{tabular}
\end{center}

\section{Notation}
\label{app:notation}

\begin{center}
\small
\begin{tabular}{@{}ll@{}}
\toprule
$n,m$ & vertices / edges of the dynamic graph\\
$\mm(G)$ & maximum matching size\\
$\RS(n,r),\ \ORS(n,r)$ & max \#\ of (ordered-)induced edge-disjoint
size-$r$ matchings (Def.~\ref{def:ors})\\
$\RS(n),\ \ORS(n)$ & the $r=\Theta(n)$ regime, as in \cite{BG24}\\
\textsc{Edges}/\textsc{Size}/\textsc{Frac}/\textsc{Witness-VC} &
maintenance objectives (\S\ref{sec:models})\\
$\OMv$ & online matrix-vector multiplication (conjecture of
\cite{HKNS15})\\
whp / adp / obl & with high probability / adaptive- / oblivious-%
adversary\\
$\Otil(\cdot)$ & suppresses $\polylog(n)$ factors\\
D1--D4, L1--L5 & sub-targets of OP-D / answer types of OP-L
(\S\ref{sec:problems})\\
\bottomrule
\end{tabular}
\end{center}

% ===============================================================
\begin{thebibliography}{99}
\small

\bibitem[ABKL23]{ABKL23}
S.~Assadi, S.~Behnezhad, S.~Khanna, and H.~Li.
On regularity lemma and barriers in streaming and dynamic matching.
In \emph{Proc.\ 55th ACM Symposium on Theory of Computing (STOC)}, 2023.
\statusM

\bibitem[ABKRS04]{ABKRS04}
N.~Alon, R.~Beigel, S.~Kasif, S.~Rudich, and B.~Sudakov.
Learning a hidden matching.
\emph{SIAM Journal on Computing}, 33(2):487--501, 2004.

\bibitem[ABR24]{ABR24}
A.~Azarmehr, S.~Behnezhad, and M.~Roghani.
Fully dynamic matching: $(2-\sqrt2)$-approximation in polylog update
time.
In \emph{Proc.\ 35th ACM-SIAM Symposium on Discrete Algorithms (SODA)},
pages 3040--3061, 2024. arXiv:2307.08772.

\bibitem[ACCSW18]{ACCSW18}
M.~Arar, S.~Chechik, S.~Cohen, C.~Stein, and D.~Wajc.
Dynamic matching: reducing integral algorithms to approximately-maximal
fractional algorithms.
In \emph{Proc.\ 45th ICALP}, 2018. \statusM

\bibitem[AKK25]{AKK25}
S.~Assadi, S.~Khanna, and P.~Kiss.
Improved bounds for fully dynamic matching via ordered
Ruzsa--Szemer\'edi graphs.
In \emph{Proc.\ 36th ACM-SIAM Symposium on Discrete Algorithms (SODA)},
pages 2971--2990, 2025. arXiv:2406.13573.

\bibitem[AKLY16]{AKLY16}
S.~Assadi, S.~Khanna, Y.~Li, and G.~Yaroslavtsev.
Maximum matchings in dynamic graph streams and the simultaneous
communication model.
In \emph{Proc.\ 27th ACM-SIAM Symposium on Discrete Algorithms (SODA)},
2016. \statusM

\bibitem[AS23]{AS23}
S.~Assadi and J.~Sundaresan.
Hidden permutations to the rescue: multi-pass semi-streaming lower
bounds for approximate matchings.
In \emph{Proc.\ 64th IEEE Symposium on Foundations of Computer Science
(FOCS)}, pages 909--932, 2023. arXiv:2310.05728.

\bibitem[BBKS25]{BBKS25}
A.~Bernstein, S.~Bhattacharya, P.~Kiss, and T.~Saranurak.
Deterministic dynamic maximal matching in sublinear update time.
In \emph{Proc.\ 57th ACM Symposium on Theory of Computing (STOC)}, 2025.
arXiv:2504.20780, DOI 10.1145/3717823.3718153.

\bibitem[BCCGNSU17]{BCCGNSU17}
J.~Blasiak, T.~Church, H.~Cohn, J.~A. Grochow, E.~Naslund, W.~F. Sawin,
and C.~Umans.
On cap sets and the group-theoretic approach to matrix multiplication.
\emph{Discrete Analysis}, 2017.

\bibitem[BCH17]{BCH17}
S.~Bhattacharya, D.~Chakrabarty, and M.~Henzinger.
Deterministic fully dynamic approximate vertex cover and fractional
matching in $O(1)$ amortized update time.
In \emph{Proc.\ 19th IPCO}, LNCS 10328, pages 86--98, 2017.

\bibitem[Beh21]{Beh21}
S.~Behnezhad.
Time-optimal sublinear algorithms for matching and vertex cover.
In \emph{Proc.\ 62nd IEEE FOCS}, pages 873--884, 2021.

\bibitem[Beh23]{Beh23}
S.~Behnezhad.
Dynamic algorithms for maximum matching size.
In \emph{Proc.\ 34th ACM-SIAM SODA}, pages 129--162, 2023.

\bibitem[Beh46]{Beh46}
F.~A. Behrend.
On sets of integers which contain no three terms in arithmetical
progression.
\emph{Proc.\ National Academy of Sciences}, 32:331--332, 1946. \statusM

\bibitem[BG24]{BG24}
S.~Behnezhad and A.~Ghafari.
Fully dynamic matching and ordered Ruzsa--Szemer\'edi graphs.
In \emph{Proc.\ 65th IEEE Symposium on Foundations of Computer Science
(FOCS)}, pages 314--327, 2024. arXiv:2404.06069.

\bibitem[BGS11]{BGS11}
S.~Baswana, M.~Gupta, and S.~Sen.
Fully dynamic maximal matching in $O(\log n)$ update time.
In \emph{Proc.\ 52nd IEEE FOCS}, pages 383--392, 2011; corrected
version: \emph{SIAM J.\ Computing} 47(3):617--650, 2018. \statusM

\bibitem[BHI15]{BHI15}
S.~Bhattacharya, M.~Henzinger, and G.~F. Italiano.
Deterministic fully dynamic data structures for vertex cover and
matching.
In \emph{Proc.\ 26th ACM-SIAM SODA}, pages 785--804, 2015.

\bibitem[BHN16]{BHN16}
S.~Bhattacharya, M.~Henzinger, and D.~Nanongkai.
New deterministic approximation algorithms for fully dynamic matching.
In \emph{Proc.\ 48th ACM STOC}, 2016. arXiv:1604.05765.

\bibitem[BK21]{BK21}
S.~Bhattacharya and P.~Kiss.
Deterministic rounding of dynamic fractional matchings.
In \emph{Proc.\ 48th ICALP}, 2021. \statusM

\bibitem[BKS23]{BKS23}
S.~Bhattacharya, P.~Kiss, and T.~Saranurak.
Dynamic $(1+\eps)$-approximate matching size in truly sublinear update
time.
In \emph{Proc.\ 64th IEEE FOCS}, 2023. arXiv:2302.05030.

\bibitem[BKSW23]{BKSW23}
S.~Bhattacharya, P.~Kiss, T.~Saranurak, and D.~Wajc.
Dynamic matching with better-than-2 approximation in polylogarithmic
update time.
In \emph{Proc.\ 34th ACM-SIAM SODA}, 2023; \emph{J.\ ACM}, DOI
10.1145/3679009. arXiv:2207.07438.

\bibitem[BLM20]{BLM20}
S.~Behnezhad, J.~\L{}\k{a}cki, and V.~Mirrokni.
Fully dynamic matching: beating 2-approximation in $\Delta^{\eps}$
update time.
In \emph{Proc.\ 31st ACM-SIAM SODA}, pages 2492--2508, 2020.

\bibitem[BRR24]{BRR24}
S.~Behnezhad, M.~Roghani, and A.~Rubinstein.
Approximating maximum matching requires almost quadratic time.
In \emph{Proc.\ 56th ACM STOC}, 2024. arXiv:2406.08595.

\bibitem[BRRS23]{BRRS23}
S.~Behnezhad, M.~Roghani, A.~Rubinstein, and A.~Saberi.
Beating greedy matching in sublinear time.
In \emph{Proc.\ 34th ACM-SIAM SODA}, 2023. arXiv:2206.13057.

\bibitem[BS16]{BS16}
A.~Bernstein and C.~Stein.
Faster fully dynamic matchings with small approximation ratios.
In \emph{Proc.\ 27th ACM-SIAM SODA}, pages 692--711, 2016.

\bibitem[CFTZ22]{CFTZ22}
M.~Christandl, O.~Fawzi, H.~Ta, and J.~Zuiddam.
Larger corner-free sets from combinatorial degenerations.
In \emph{Proc.\ 13th ITCS}, 2022. \statusM

\bibitem[CGL15]{CGL15}
R.~Clifford, A.~Gr{\o}nlund, and K.~G. Larsen.
New unconditional hardness results for dynamic and online problems.
In \emph{Proc.\ 56th IEEE FOCS}, 2015. \statusM

\bibitem[FLNRRS02]{FLNRRS02}
E.~Fischer, E.~Lehman, I.~Newman, S.~Raskhodnikova, R.~Rubinfeld, and
A.~Samorodnitsky.
Monotonicity testing over general poset domains.
In \emph{Proc.\ 34th ACM STOC}, 2002. \statusM\ (attribution of the
$\ORS(\Theta(n))\ge n^{1/\Theta(\log\log n)}$ construction per
\cite{BG24}.)

\bibitem[GP13]{GP13}
M.~Gupta and R.~Peng.
Fully dynamic $(1+\eps)$-approximate matchings.
In \emph{Proc.\ 54th IEEE FOCS}, pages 548--557, 2013.

\bibitem[GSSU22]{GSSU22}
F.~Grandoni, C.~Schwiegelshohn, S.~Solomon, and A.~Uzrad.
Maintaining an EDCS in general graphs: simpler, density-sensitive and
with worst-case time bounds.
In \emph{Proc.\ 5th SIAM SOSA}, pages 12--23, 2022. \statusM

\bibitem[HKNS15]{HKNS15}
M.~Henzinger, S.~Krinninger, D.~Nanongkai, and T.~Saranurak.
Unifying and strengthening hardness for dynamic problems via the online
matrix-vector multiplication conjecture.
In \emph{Proc.\ 47th ACM STOC}, pages 21--30, 2015.

\bibitem[Kis22]{Kis22}
P.~Kiss.
Improving update times of dynamic matching algorithms from amortized to
worst case.
In \emph{Proc.\ 13th ITCS}, 94:1--94:21, 2022.

\bibitem[Lar12]{Lar12}
K.~G. Larsen.
The cell probe complexity of dynamic range counting.
In \emph{Proc.\ 44th ACM STOC}, 2012. \statusM

\bibitem[Liu24]{Liu24}
Y.~P. Liu.
On approximate fully-dynamic matching and online matrix-vector
multiplication.
In \emph{Proc.\ 65th IEEE FOCS}, pages 228--243, 2024. arXiv:2403.02582.
(General-graph extension: arXiv:2503.01147.)

\bibitem[LW17]{LW17}
K.~G. Larsen and R.~R. Williams.
Faster online matrix-vector multiplication.
In \emph{Proc.\ 28th ACM-SIAM SODA}, pages 2182--2189, 2017.

\bibitem[NS13]{NS13}
O.~Neiman and S.~Solomon.
Simple deterministic algorithms for fully dynamic maximal matching.
In \emph{Proc.\ 45th ACM STOC}, 2013. arXiv:1207.1277.

\bibitem[OR10]{OR10}
K.~Onak and R.~Rubinfeld.
Maintaining a large matching and a small vertex cover.
In \emph{Proc.\ 42nd ACM STOC}, 2010. \statusM

\bibitem[Pat10]{Pat10}
M.~P\u{a}tra\c{s}cu.
Towards polynomial lower bounds for dynamic problems.
In \emph{Proc.\ 42nd ACM STOC}, 2010. \statusM

\bibitem[Pra26]{Pra26}
K.~Pratt.
A note on ordered Ruzsa--Szemer\'edi graphs.
In \emph{Proc.\ 9th SIAM Symposium on Simplicity in Algorithms (SOSA)},
2026. arXiv:2502.02455, DOI 10.1137/1.9781611978964.27.

\bibitem[RS78]{RS78}
I.~Z. Ruzsa and E.~Szemer\'edi.
Triple systems with no six points carrying three triangles.
In \emph{Combinatorics (Keszthely, 1976)}, Coll.\ Math.\ Soc.\ J.\
Bolyai 18, pages 939--945, 1978.

\bibitem[Sol16]{Sol16}
S.~Solomon.
Fully dynamic maximal matching in constant update time.
In \emph{Proc.\ 57th IEEE FOCS}, pages 325--334, 2016.

\bibitem[Waj20]{Waj20}
D.~Wajc.
Rounding dynamic matchings against an adaptive adversary.
In \emph{Proc.\ 52nd ACM STOC}, 2020. \statusM

\bibitem{AS19} S.~Assadi, S.~Solomon. When algorithms for maximal
independent set and maximal matching run in sublinear time. ICALP
2019, 17:1--17:17. (Thm.~3: deterministic $\Omega(n^2)$ for finding a
maximal matching at neighborhood independence $\beta=2$.)
\bibitem{ACK19} S.~Assadi, Y.~Chen, S.~Khanna. Sublinear algorithms
for $(\Delta+1)$ vertex coloring. SODA 2019, pp.~767--786.
(Thm.~6: finding a maximal matching needs $\Omega(n^2)$ queries,
randomized.)
\end{thebibliography}

\end{document}
